% SGA 3, Expose V, Loop 1 English reconstruction.
% Scope: Section 6 reduction, Lemma 6.1, and the opening of proof (a).
% Authority: Polo--Gille Exp5-13oct24.pdf, component-local p. 21
% (combined-reader p. 283), re-edited page header 271.
% Jacob Reinhold's CC BY 4.0 Markdown is comparison/drafting material only.

这是 \(U_\ast\) 在基变换 \(u:V\to U\) 下的逆像,
因而也是 \(X_\ast\) 在基变换
\SGAThreeDiagram[0.43\textwidth]
  {figures/exp5/p271_section6_base_change_d0.png}
  {定义 \(V_\ast\) 的基变换.}
  {fig:sga3-V-6-base-change-d0}
  {第 V 篇, 分篇第~21 页 / 合订本第~283 页 / 重版本页眉页码~271}
下的逆像.

由 \hyperref[par:sga3-V-3-c]{\S3.c)}, \(V_\ast\)
同构于群胚 \(V'_\ast\), 后者是 \(X_\ast\) 在基变换
\SGAThreeDiagram[0.49\textwidth]
  {figures/exp5/p271_section6_base_change_d1.png}
  {定义 \(V'_\ast\) 的基变换.}
  {fig:sga3-V-6-base-change-d1}
  {第 V 篇, 分篇第~21 页 / 合订本第~283 页 / 重版本页眉页码~271}
下的逆像. 因此 \(V'_\ast\) 在 \((\mathbf{Sch}/S)\) 中容许余核.
由于 \(v:V\to X_0\) 平坦、满射且有限,
它忠实平坦且拟紧, 因此由第 IV 篇, 6.3.2, 它是普遍有效满态射.
\footnote{法文重版本误引为第 III 篇, 该篇止于 \S4;
见 Canon 记录 \texttt{SGA-SRC-000857}.}
因而命题~\hyperref[prop:sga3-V-3.1]{3.1} 表明,
\(X_\ast\) 在 \((\mathbf{Sch}/S)\) 中容许余核
\(\operatorname{Coker}(d_0,d_1)\).
至此证明了\hyperref[lem:sga3-V-6.1]{下面引理}的
\hyperref[lem:sga3-V-6.1-i]{(i) 项}的第一个断言.
\footnote{\label{note:sga3-V-32}编者注: 下面的论证略有修改.
特别地, \(d_0\) 平坦这一附加假设已移到
\hyperref[lem:sga3-V-6.1-iv]{(iv) 项}; 原来的 (ii) 项已分成
\hyperref[lem:sga3-V-6.1-ip]{(i') 项}和~\hyperref[lem:sga3-V-6.1-ii]{(ii) 项};
并新增了\hyperref[lem:sga3-V-6.1-ip]{(i') 项}的第二个断言.}

\medskip
\phantomsection\label{lem:sga3-V-6.1}%
\noindent\textbf{引理 6.1.}---\emph{假设
\((\mathbf{Sch}/S)\)-群胚 \(X_\ast\) 具有拟截面. 则:}
\begin{enumerate}
\item[\SGAInlineRefTarget{lem:sga3-V-6.1:item:i}(i)]\phantomsection\label{lem:sga3-V-6.1-i}%
\emph{\((d_0,d_1)\) 在 \((\mathbf{Sch}/S)\) 中有余核 \((Y,p)\).
  此外, 每个这样的 \((Y,p)\) 都是所有赋环空间范畴中
  \((d_0,d_1)\) 的余核.}
\item[(i')]\phantomsection\label{lem:sga3-V-6.1-ip}%
\emph{态射 \(p\) 满射; 若 \(d_0\) 开
  (分别地, 普遍闭), 则它也开 (分别地, 普遍闭).}
\item[(ii)]\phantomsection\label{lem:sga3-V-6.1-ii}%
\emph{假设 \(S\) 局部 Noether, 且 \(X_0\) 在 \(S\) 上
  局部有限型 (分别地, 有限型). 则 \(p\) 和 \(Y\to S\)
  都局部有限呈示 (分别地, 有限呈示).}
\item[(iii)]\phantomsection\label{lem:sga3-V-6.1-iii}%
\emph{以 \(d_0\) 和 \(d_1\) 为分量的态射
  \(X_1\to X_0\times_Y X_0\) 满射.}
\item[(iv)]\phantomsection\label{lem:sga3-V-6.1-iv}%
\emph{若 \((d_0,d_1)\) 是等价对, 则
  \(X_1\to X_0\times_Y X_0\) 是同构.
  此外, 若 \(d_0:X_1\to X_0\) 平坦, 则 \(p\) 忠实平坦.}
\end{enumerate}

在证明\hyperref[lem:sga3-V-6.1-i]{(i) 项}的第二个断言之前,
先证明 \hyperref[lem:sga3-V-6.1-ip]{(i') 项},
\hyperref[lem:sga3-V-6.1-ii]{(ii) 项}及 \hyperref[lem:sga3-V-6.1-iii]{(iii) 项}.

\medskip
\phantomsection\label{par:sga3-V-6-a}%
\noindent\textbf{a) \hyperref[lem:sga3-V-6.1-ip]{(i')} 与 \hyperref[lem:sga3-V-6.1-ii]{(ii)} 的证明.}---%
刚才已见 \((Y,p)\) 既可与
\(\operatorname{Coker}(v_0,v_1)\) 等同, 也可与
\(\operatorname{Coker}(u_0,u_1)\) 等同.
设 \(q\) 和 \(r\) 分别是从 \(U\) 和 \(V\) 到 \(Y\) 的典范满态射:
\SGAThreeDiagram[0.39\textwidth]
  {figures/exp5/p271_lemma61_cokernel_triangle.png}
  {到 \(Y\) 的典范满态射.}
  {fig:sga3-V-6.1-cokernel-triangle}
  {第 V 篇, 分篇第~21 页 / 合订本第~283 页 / 重版本页眉页码~271}
由假设, \(v\) 满射且有限局部自由, 因而开.

% Translation coverage: entire assigned component; adjoining components retain the original statement/proof seams.
% Source lines: 84; translated body lines before this audit comment: 80.
% Math preserved: 46 inline formulas, 0 display blocks, 0 equation environments; mathematical prose translated where present.
% Labels: 8/8; hyperrefs: 14/14; SGARefLinks: 0/0; inline targets: 1/1; footnotes: 2/2.
% Figure arguments preserved; all rendered captions and natural-language source locators translated.
% Unresolved TODOs: 0; compilation and visual inspection remain with the volume integrator.
